\documentclass{article}
\usepackage[utf8]{inputenc}
\usepackage[T1]{fontenc}
\usepackage[polish]{babel}
\usepackage{amsmath}
\usepackage{amssymb}
\usepackage{graphicx}
\usepackage{hyperref}
\usepackage{listings}
\usepackage{xcolor}

% Code highlighting settings
\definecolor{codegray}{rgb}{0.5,0.5,0.5}
\definecolor{codebackground}{rgb}{0.95,0.95,0.95}
\lstset{
    backgroundcolor=\color{codebackground},
    commentstyle=\color{codegray},
    keywordstyle=\color{blue},
    numberstyle=\tiny\color{codegray},
    stringstyle=\color{teal},
    basicstyle=\footnotesize\ttfamily,
    breakatwhitespace=false,
    breaklines=true,
    captionpos=b,
    keepspaces=true,
    numbers=left,
    numbersep=5pt,
    showspaces=false,
    showstringspaces=false,
    showtabs=false,
    tabsize=2
}

\title{Dokumentacja Projektu: System RAG do Programu Studiów}
\author{Stefan Nosek}
\date{\today}

\begin{document}

\maketitle

\tableofcontents
\newpage

\section{Wstęp}
Sprawozdanie przedstawia dokumentację projektu ``System RAG do programu studiów''. Projekt ten miał na celu stworzenie aplikacji webowej umożliwiającej semantyczne wyszukiwanie i odpowiadanie na pytania dotyczące programu studiów Politechniki Łódzkiej, wykorzystując technikę Retrieval-Augmented Generation (RAG). System ten ma za zadanie ułatwić użytkownikom dostęp do informacji o kursach, wymaganiach, efektach uczenia się oraz innych szczegółach związanych z ofertą edukacyjną uczelni.

\section{Cel projektu}
Głównym celem projektu jest rozwiązanie problemu trudności w szybkim i efektywnym wyszukiwaniu szczegółowych informacji w bazie danych programów studiów Politechniki Łódzkiej. Aplikacja ma dostarczyć intuicyjny interfejs użytkownika, który na podstawie naturalnego języka zapytań będzie w stanie odnaleźć najbardziej relewantne kursy i wygenerować zwięzłe odpowiedzi. Cele szczegółowe to:
\begin{itemize}
    \item Zaprojektowanie i implementacja modułu do scrapowania danych z witryn internetowych Politechniki Łódzkiej zawierających informacje o programach studiów.
    \item Budowa i utrzymanie bazy danych przechowującej przetworzone dane o kursach wraz z ich wektorowymi reprezentacjami (embeddingami).
    \item Opracowanie silnika wyszukiwania semantycznego, zdolnego do odnajdywania najbardziej podobnych kursów do zadanego zapytania.
    \item Integracja z modelem językowym (LLM) w celu generowania spójnych i merytorycznych odpowiedzi na podstawie wyszukanych informacji.
    \item Stworzenie interfejsu użytkownika (frontendu) umożliwiającego interakcję z systemem.
\end{itemize}

\section{Przegląd istniejących rozwiązań}
W tej sekcji należy przeprowadzić analizę trzech podobnych rozwiązań dostępnych na rynku. Mogą to być systemy informacyjne uczelni, wyszukiwarki kursów, czy też inne aplikacje wykorzystujące techniki RAG lub semantycznego wyszukiwania w kontekście edukacyjnym. Dla każdego rozwiązania należy opisać jego funkcjonalności, mocne strony, ograniczenia oraz technologie, w jakich zostało zrealizowane.
\textit{[Wstaw tutaj analizę istniejących rozwiązań]}

\section{Opis architektury}
System został zaprojektowany jako aplikacja webowa w architekturze mikrousługowej, składającej się z kilku głównych komponentów współpracujących ze sobą. Poniżej przedstawiono szczegółowy opis struktury systemu, stosowanych technologii i sposobu komunikacji między komponentami.

\subsection{Komponenty systemu}
\begin{itemize}
    \item \textbf{Scraper danych}: Moduł odpowiedzialny za zbieranie danych o kursach z oficjalnych stron Politechniki Łódzkiej.
    \item \textbf{Baza danych PostgreSQL z rozszerzeniem PGVector}: Przechowuje dane o kursach w formacie JSONB oraz ich wektorowe reprezentacje (embeddingi). Dodatkowo umożliwia efektywne wyszukiwanie wektorowe.
    \item \textbf{Silnik wyszukiwania semantycznego}: Odpowiada za generowanie embeddingów dla zapytań użytkowników oraz wyszukiwanie najbardziej podobnych wektorów w bazie danych.
    \item \textbf{Moduł chatbota}: Wykorzystuje duży model językowy (LLM), dostępny poprzez darmowe API, do generowania odpowiedzi na pytania użytkowników w oparciu o kontekst dostarczony przez silnik wyszukiwania semantycznego.
    \item \textbf{Aplikacja webowa (Streamlit)}: Stanowi interfejs użytkownika, za pomocą którego użytkownicy mogą zadawać pytania i otrzymywać odpowiedzi.
    \item \textbf{Docker i Docker Compose}: Używane do konteneryzacji i orkiestracji wszystkich komponentów systemu.
\end{itemize}

\subsection{Stos technologiczny}
\begin{itemize}
    \item \textbf{Python}: Główny język programowania użyty do implementacji scrapera, logiki backendu, silnika semantycznego wyszukiwania i chatbota.
    \item \textbf{Streamlit}: Framework do budowy interfejsu użytkownika webowego.
    \item \textbf{PostgreSQL z PGVector}: System zarządzania bazą danych z rozszerzeniem do obsługi embeddingów.
    \item \textbf{SQLAlchemy}: Używany do interakcji z bazą danych z poziomu Pythona (kwestie ORM pominięte).
    \item \textbf{sentence-transformers}: Biblioteka do generowania embeddingów.
    \item \textbf{Groq API}: Wykorzystywane do dostępu do modeli językowych (LLM).
    \item \textbf{Selenium}: Biblioteka do automatyzacji przeglądarek, używana w scraperze.
    \item \textbf{Pandas, NumPy, python-dotenv}: Dodatkowe biblioteki wspomagające przetwarzanie danych i działanie aplikacji.
\end{itemize}

\section{Opis implementacji}
Implementacja projektu została podzielona na trzy główne obszary: backend, bazę danych i frontend.

\subsection{Backend i logika aplikacji}
Logika backendu obejmuje scraper, moduły do obsługi bazy danych (populacja i generowanie embeddingów), silnik semantycznego wyszukiwania oraz moduł chatbota.

\subsubsection{Moduły scrapera (\texttt{scraper/})}
Scraper składa się z trzech głównych skryptów:
\begin{itemize}
    \item \texttt{helper.py}: Zawiera funkcję \texttt{extract\_link} do parsowania linków z atrybutów \texttt{onclick}.
    \item \texttt{parse\_top\_level.py}: Odpowiedzialny za pobieranie linków do programów studiów, kierunków, stopni i form studiów z głównej strony (\url{https://programy.p.lodz.pl/ectslabel-web/}). Wyniki są zapisywane do pliku \texttt{top\_level.csv}. Używa Selenium do nawigacji po stronie i Pandas do zarządzania danymi.
    \item \texttt{parse\_mid\_level.py}: Przetwarza linki z \texttt{top\_level.csv} w celu ekstrakcji kodów kursów i linków do szczegółowych kart kursów. Radzi sobie z różnymi specjalizacjami i semestrami, zapisując dane do \texttt{mid\_level.csv}. Obsługuje również przypadki "Mobility Semester".
    \item \texttt{parse\_low\_level.py}: Pobiera szczegółowe dane z indywidualnych kart przedmiotów. Rozróżnia standardową i alternatywną strukturę kart. Dane są zapisywane jako pliki JSON dla każdego kursu w katalogu \texttt{data/dane\_low\_level}. Wykorzystuje Selenium do interakcji ze stroną i obsługuje dynamiczne ładowanie elementów.
\end{itemize}

\subsubsection{Moduły bazy danych (\texttt{db/})}
\begin{itemize}
    \item \texttt{populate\_db.py}: Skrypt do inicjalizacji schematu bazy danych i populacji jej danymi z plików JSON zebranych przez scraper. Tworzy tabelę \texttt{course\_data} z kolumnami \texttt{id}, \texttt{data} (JSONB) i \texttt{embedding} (VECTOR(384)).
    \item \texttt{generate\_embeddings.py}: Skrypt odpowiedzialny za generowanie embeddingów dla opisów kursów. Wykorzystuje model \texttt{sentence-transformers/\allowbreak paraphrase-multilingual-MiniLM-L12-v2} do kodowania opisów kursów, które są konstruowane z różnych pól danych kursu. Następnie aktualizuje pole \texttt{embedding} w tabeli \texttt{course\_data}.
\end{itemize}

\subsubsection{Moduły RAG (\texttt{src/rag/})}
\begin{itemize}
    \item \texttt{db\_driver.py}: Zawiera funkcje do łączenia się z bazą danych PostgreSQL za pomocą SQLAlchemy i \texttt{psycopg2-binary}. Używa zmiennych środowiskowych do konfiguracji połączenia i udostępnia context managera \texttt{get\_db} do bezpiecznej pracy z sesjami bazy danych.
    \item \texttt{semantic\_search.py}: Implementuje klasę \texttt{SemanticSearchEngine}, która wykorzystuje model \texttt{sentence-transformers/paraphrase-multilingual\allowbreak MiniLM-L12-v2} do generowania embeddingów zapytań. Wykonuje zapytanie do bazy danych, wykorzystując operator \texttt{<=>} (odległość cosinusowa) do znalezienia 10 najbliższych wektorów kursów.
    \item \texttt{chatbot.py}: Zawiera funkcje \texttt{create\_prompt} i \texttt{get\_answer}. Funkcja \\\texttt{create\_prompt} formatuje zapytanie użytkownika i informacje o najbliższych kursach w spójny prompt dla modelu językowego. Funkcja \texttt{get\_answer} wysyła sformułowany prompt do API Groq (\texttt{llama-3.3-70b-versatile}) i zwraca wygenerowaną odpowiedź.
\end{itemize}

\subsection{Baza danych}
Baza danych opiera się na PostgreSQL z rozszerzeniem PGVector, co umożliwia przechowywanie i efektywne wyszukiwanie wektorów o wysokiej wymiarowości.
\begin{itemize}
    \item \textbf{Tabela \texttt{course\_data}}: Przechowuje surowe dane kursów w kolumnie \texttt{data} typu JSONB oraz ich embeddingi w kolumnie \texttt{embedding} typu VECTOR(384).
    \item \textbf{Konfiguracja}: Konfiguracja połączenia z bazą danych odbywa się za pomocą zmiennych środowiskowych (\texttt{POSTGRES\_USER}, \texttt{POSTGRES\_PASSWORD}, \texttt{POSTGRES\_DB}, \texttt{POSTGRES\_PORT}).
\end{itemize}

\subsection{Frontend}
Frontend aplikacji został zaimplementowany w Streamlit.
\begin{itemize}
    \item \texttt{webapp/app.py}: Główny plik aplikacji Streamlit. Użytkownik wprowadza zapytanie tekstowe. Po naciśnięciu przycisku ``Pytaj'' aplikacja wykorzystuje \texttt{SemanticSearchEngine} do wyszukania podobnych kursów. Następnie generuje prompt i wysyła go do chatbota, wyświetlając otrzymaną odpowiedź.
    \item \texttt{webapp/Dockerfile}: Definiuje środowisko dla aplikacji Streamlit, bazując na obrazie \texttt{python:3.9-slim}. Instaluje niezbędne zależności systemowe i Pythonowe z pliku \texttt{requirements.txt}. Kopiuje kod aplikacji i ustawia punkt wejścia do uruchomienia aplikacji Streamlit.
\end{itemize}

\subsection{Konteneryzacja (Docker i Docker Compose)}
Cała aplikacja jest konteneryzowana przy użyciu Docker. Plik \texttt{docker-compose.yml} definiuje dwa serwisy:
\begin{itemize}
    \item \textbf{postgres}: Wykorzystuje obraz \texttt{ankane/pgvector:latest} do uruchomienia bazy danych PostgreSQL z rozszerzeniem PGVector. Wszystko konfigurowane jest z pomocą zmiennych środowiskowych.
    \item \textbf{streamlit}: Budowany z \texttt{webapp/Dockerfile}. Mapuje port 8501 na hosta i montuje woluminy \texttt{./webapp} oraz \texttt{./src} do kontenera. Przekazuje zmienne środowiskowe potrzebne do połączenia z bazą danych i klucza API Groq. Posiada zależność od serwisu \texttt{postgres}.
\end{itemize}
Ten układ zapewnia izolację środowisk, łatwość wdrożenia i odtwarzalność.

\section{Podsumowanie i wnioski}
Projekt ``System RAG do programu studiów'' stanowi kompleksowe rozwiązanie problemu dostępu do informacji o kursach na Politechnice Łódzkiej. Dzięki zastosowaniu technik RAG, system jest w stanie nie tylko odnajdywać relewantne dokumenty, ale także generować spójne i kontekstowe odpowiedzi, znacząco poprawiając doświadczenie użytkownika w porównaniu do tradycyjnych wyszukiwarek.

\subsection{Napotkane wyzwania}
Podczas realizacji projektu napotkano kilka kluczowych wyzwań. Jednym z nich była kwestia doboru modelu dla silnika wyszukiwania semantycznego oraz chatbota. Wybór odpowiedniego modelu Sentence-Transformers do generowania embeddingów oraz modelu językowego (LLM) z oferty Groq wymagał empirycznego badania jakości poprzez ręczne testy i weryfikację trafności wyników. Zdarzało się, że pierwsze wybrane modele nie dawały satysfakcjonujących rezultatów, co wymuszało eksperymentowanie z innymi opcjami i iteracyjne doskonalenie.

Kolejnym istotnym wyzwaniem było kilkupoziomowe scrapowanie danych. Strony Politechniki Łódzkiej charakteryzują się złożoną strukturą, w której informacje o kursach są zagnieżdżone na różnych poziomach (programy -> kierunki -> specjalizacje -> semestry -> kursy -> szczegóły kursu). Dodatkowo, napotkano na problemy z pustymi elementami, brakującymi danymi oraz różnym formatowaniem kart przedmiotów, co wymagało implementacji elastycznych mechanizmów parsowania i obsługi błędów w skryptach scrapera. Szczególnym przypadkiem był ``Mobility Semester'', który wymagał specjalnego traktowania ze względu na brak szczegółowych informacji o kursach.

Trzecim znaczącym problemem było nawiązanie połączenia z bazą danych PostgreSQL uruchomioną w kontenerze Docker. Początkowo połączenie nie działało, ponieważ port 5432, domyślny dla PostgreSQL, był już zajęty przez lokalnie zainstalowaną wersję PostgreSQL. Problem został rozwiązany po zbadaniu otwartych portów i zmianie mapowania portów w pliku \texttt{docker-compose.yml} na wolny port hosta.

\subsection{Możliwości dalszego rozwoju}
\begin{itemize}
    \item Rozszerzenie zakresu danych o kolejne programy studiów, lata akademickie, czy inne źródła informacji (np. regulaminy, opisy katedr).
    \item Wdrożenie bardziej zaawansowanych technik RAG, takich jak re-ranking wyników, multistep reasoning, czy wykorzystanie hybrydowego wyszukiwania (semantyczne + słowa kluczowe).
    \item Usprawnienie interfejsu użytkownika, dodanie funkcji filtrowania, sortowania wyników, czy wizualizacji danych.
    \item Optymalizacja wydajności silnika semantycznego wyszukiwania i chatbota dla dużych wolumenów danych.
    \item Dodanie mechanizmu feedbacku od użytkowników w celu doskonalenia jakości odpowiedzi.
    \item Implementacja uwierzytelniania i autoryzacji dla różnych typów użytkowników.
\end{itemize}

\section{Bibliografia}
\begin{thebibliography}{99}
    \bibitem{docker} Dokumentacja Docker: \url{https://docs.docker.com/}
    \bibitem{docker-compose-docs} Dokumentacja Docker Compose: \url{https://docs.docker.com/compose/}
    \bibitem{streamlit-docs} Dokumentacja Streamlit: \url{https://docs.streamlit.io/}
    \bibitem{postgres-docs} Dokumentacja PostgreSQL: \url{https://www.postgresql.org/docs/}
    \bibitem{pgvector-github} Repozytorium GitHub PGVector: \url{https://github.com/pgvector/pgvector}
    \bibitem{python-docs} Dokumentacja Python: \url{https://docs.python.org/3/}
    \bibitem{sqlalchemy-docs} Dokumentacja SQLAlchemy: \url{https://docs.sqlalchemy.org/}
    \bibitem{sentence-transformers-docs} Dokumentacja Sentence-Transformers: \url{https://www.sbert.net/docs/index.html}
    \bibitem{groq-docs} Dokumentacja Groq API: \url{https://console.groq.com/docs/api-reference}
    \bibitem{selenium-docs} Dokumentacja Selenium: \url{https://www.selenium.dev/documentation/webdriver/}
    \bibitem{pandas-docs} Dokumentacja Pandas: \url{https://pandas.pydata.org/docs/}
    \bibitem{numpy-docs} Dokumentacja NumPy: \url{https://numpy.org/doc/}
\end{thebibliography}

\end{document}